\addtocontents{toc}{\protect\Needspace{3\baselineskip}}
\originalchapter{Wilson 环与热全息}\label{ch:Wilson}
\originalsection{静态重夸克势}
边界矩形环时间长度 $\mathcal T\gg L$ 时 $\langle W\rangle\sim\ee^{-\mathcal TV(L)}$. 在经典弦近似, 以终止在边界环上的世界面最小面积估计期望值. AdS$_5$ 欧氏度规中取静态规范 $(t,x)$、轮廓 $z(x)$, Nambu--Goto 每时间作用量为
\begin{equation}
\frac S{\mathcal T}=\frac{R^2}{2\pi\alpha'}\int_{-L/2}^{L/2}
\frac{\sqrt{1+z'^2}}{z^2}\dd x.
\end{equation}
因 Lagrangian 无显式 $x$, $1/[z^2\sqrt{1+z'^2}]=1/z_*^2$, 得 $z'^2=z_*^4/z^4-1$. 因而 $L=2z_*\int_0^1y^2\dd y/\sqrt{1-y^4}$. 两条直弦的无限端点质量须减去, 才得有限势
$V(L)=-4\pi^2\sqrt\lambda/[\Gamma(1/4)^4L]$.
这个结果适用带标量耦合的 $\mathcal N=4$ Wilson 算符、强耦合与大 $N$. 系数的 $\sqrt\lambda$ 与弱耦合的 $\lambda$ 不同; 两者不处在相同的微扰区间.

\originalsection{黑膜与有限温度}
有限温度态对应平直 AdS 黑膜, $\dd s^2=R^2z^{-2}[-f(z)\dd t^2+\dd\boldsymbol x^2+\dd z^2/f(z)]$, $f=1-(z/z_h)^d$. 欧氏正则性给 $T=d/(4\pi z_h)$. 视界面积每坐标空间体积给熵密度 $s=R^{d-1}/(4G_{d+1}z_h^{d-1})$, 与边界 CFT 的 $T^{d-1}$ 量纲一致.

Wilson 世界面可以是连接的 U 形面, 也可由两条各到视界的直弦组成. 比较\textbf{同一减法方案}下的自由能, 距离大时断开构型占优, 给重夸克颜色屏蔽. 连接解存在的最大距离与实际自由能交叉距离不同; 不能只看到连接解消失便把该点叫热力学转变点. 具体积分和数值由入选 PS5.3 展开.

加入体规范场可描述化学势, 其边界值对应源, 内部正则性要求以视界为参照的势差. 本卷未系统求带电黑膜、输运系数和非平衡响应, 将这些原课程高阶范围明确留作后续扩编.

\originalsection{零温势系数的积分}
长度积分令 $u=y^4$, 得 $I=\int_0^1y^2\dd y/\sqrt{1-y^4}=\frac14B(3/4,1/2)=\sqrt\pi\Gamma(3/4)/\Gamma(1/4)$. 所以 $L=2Iz_*$. 减去两条到 Poincaré 深处的直弦后, 有
\[
V=\frac{R^2}{\pi\alpha'z_*}
\left[\int_0^1\left(\frac1{y^2\sqrt{1-y^4}}-\frac1{y^2}\right)\dd y-1\right]
=-\frac{R^2I}{\pi\alpha'z_*}.
\]
最后等式可通过把 $\sqrt{1-y^4}/y$ 的导数与积分组合复核, 端点的减法正好消掉 $y\to0$ 的发散. 消去 $z_*$ 得 $V=-2I^2\sqrt\lambda/(\pi L)$. 用 $\Gamma(1/4)\Gamma(3/4)=\pi\sqrt2$ 得正文系数 $4\pi^2/\Gamma(1/4)^4$. 若只求最小面而不减端点自能, 并不能得到一个可比较的有限夸克势.
